
\medskip
\phantomsection\label{prop:sga3-VIB-7.9}%
\noindent\textbf{Proposition 7.9.---}%
\emph{Let \(k\) be an algebraically closed field, and let \(G\)
be a \(k\)-group locally of finite type.}
\begin{enumerate}
\item[\textup{\SGAInlineRefTarget{prop:sga3-VIB-7.9:item:i}(i)}]\phantomsection\label{concl:sga3-VIB-7.9-i}
  \emph{Let \(A\) and \(B\) be two ind-constructible subsets
  of \(G\). Denote by \(A_0\) the set of rational points of
  \(G\) that belong to \(A\). Then}
  \[
    (A\cdot B)_0=A_0\cdot B_0,
  \]
  \emph{where the product on the right is taken in the group
  \(G(k)\).}

\item[\textup{\SGAInlineRefTarget{prop:sga3-VIB-7.9:item:ii}(ii)}]\phantomsection\label{concl:sga3-VIB-7.9-ii}
  \emph{Let \(X\) be a geometrically reduced \(k\)-scheme
  locally of finite type, and let \(f:X\to G\) be a
  \(k\)-morphism. Put}
  \[
    \Gamma'_G(f)=
      \bigcup_{n\geq1}f^n(X^n).
  \]
  \emph{Then \(\Gamma'_G(f)_0\) is the subgroup of \(G(k)\)
  generated by \(f(X)_0\).}

\item[\textup{(iii)}]\phantomsection\label{concl:sga3-VIB-7.9-iii}
  \emph{In particular, let \(A\) and \(B\) be two smooth
  group subschemes of \(G\), and put}
  \[
    \Gamma'=
      \bigcup_{n\geq1}
        \nu^n\bigl((A\times_k B)^n\bigr)
  \]
  \emph{with the notation of
  \hyperref[def:sga3-VIB-7.2.2]{7.2.2}. Then
  \(\Gamma'_0\) is the commutator subgroup of \(A(k)\) and
  \(B(k)\) in \(G(k)\).}
\end{enumerate}

\emph{Proof.}
We first prove \textup{\SGARefLink{prop:sga3-VIB-7.9:item:i}{(i)}}. It is clear that
\[
  A_0\cdot B_0\subset(A\cdot B)_0.
\]
Conversely, let \(z\in(A\cdot B)_0\). Then
\(\mu^{-1}(z)\) is a closed subset of \(G\times_kG\), while
\(A\times_kB\) is an ind-constructible subset of
\(G\times_kG\) (cf.\ \hyperref[par:sga3-VIB-3.0]{3.0}).
Consequently,
\[
  \mu^{-1}(z)\cap(A\times_kB)
\]
is a nonempty ind-constructible subset of \(G\times_kG\).
By EGA~IV\(_3\), 10.4.8, it therefore contains a rational
point of \(G\times_kG\). Its projections \(x\) and \(y\) are
rational points of \(G\) such that \(x\in A\), \(y\in B\),
and \(x\cdot y=z\). Hence
\begingroup
\renewcommand{\thefootnote}{\fnsymbol{footnote}}%
\footnote[1]{\label{note:sga3-VIB-7.9-product}Translator's note:
the French proof prints \((A,B)_0=A_0\cdot B_0\) here.
The statement and the multiplication argument require
\((A\cdot B)_0=A_0\cdot B_0\), which is restored in the
English text.}%
\endgroup
\[
  (A\cdot B)_0=A_0\cdot B_0.
\]

To prove \textup{\SGARefLink{concl:sga3-VIB-7.9-ii}{(ii)}}, observe that, since \(f^n\) is locally
of finite type, \(f^n(X^n)\) is an ind-constructible subset
of \(G\) (EGA~IV\(_1\), 1.9.5\textup{(viii)})\begingroup
\renewcommand{\thefootnote}{\fnsymbol{footnote}}%
\footnote[2]{\label{note:sga3-VIB-7.9-EGA-volume}
Translator's note: the French source prints EGA~IV\(_4\) here.
The cited result 1.9.5\textup{(viii)} belongs to
EGA~IV\(_1\), which is restored in the English text.}%
\endgroup.
Assertion \textup{\SGARefLink{prop:sga3-VIB-7.9:item:i}{(i)}} then shows by induction that, if
\[
  A=f(X)_0,
\]
one has
\[
  f^n(X^n)_0=(A\cup A^{-1})^n.
\]
Consequently,
\phantomsection\label{form:sga3-VIB-7.9-generated-rational-points}%
\[
  \Gamma'_G(f)_0
    =\bigcup_{n\geq1}f^n(X^n)_0
    =\bigcup_{n\geq1}(A\cup A^{-1})^n,
\]
which is the subgroup of \(G(k)\) generated by
\(A=f(X)_0\).

Finally, \textup{(iii)} follows from \textup{\SGARefLink{prop:sga3-VIB-7.9:item:ii}{(ii)}} and the
definitions.

\medskip
\phantomsection\label{cor:sga3-VIB-7.10}%
\noindent\textbf{Corollary 7.10.---}%
\emph{Under the hypotheses of
\hyperref[prop:sga3-VIB-7.8]{7.8}, if \(k\) is algebraically
closed, then \((A,B)(k)\) is the commutator subgroup of
\(A(k)\) and \(B(k)\) in \(G(k)\).}

Indeed, it is enough to apply
\hyperref[concl:sga3-VIB-7.9-iii]{7.9\textup{(iii)}},
\hyperref[prop:sga3-VIB-7.8]{7.8}, and
\hyperref[prop:sga3-VIB-7.6]{7.6}.
